\section{Остановка с гарантией точности и подбором шага}
\label{sec:stopping}
Предположим, что выполнены условия выпуклой модели из
раздела~\ref{sec:model}. Из выпуклости следует
\begin{equation}\label{eq:gap_certificate}
 f(x)-f(x^*)\le
 \langle\nabla f(x),x-x^*\rangle\le G(x).
\end{equation}
Полный зазор даёт вычисляемую границу ошибки и при
$\nabla f(x^*)\ne0$. Оценивать этот неизвестный градиент не нужно.
В доказательстве используем положительную границу кривизны $K_Q$
из \eqref{eq:global_curvature}. Если она равна нулю, её можно
увеличить до любого положительного числа. Сам алгоритм не требует
заданных значений $K_Q$, $L_0$ или $L_1$.

\paragraph{Алгоритм с точными оракулами.}
Зададим $\varepsilon>0$ и начальное значение кривизны
$L_{\rm init}>0$. В текущей точке $x$ вычисляем градиент,
линейный минимум $s\in Q$ и зазор $G(x)$. Сохраняем
\[
 \ell\leftarrow\max\{\ell,f(x)-G(x)\},
\]
причём первый запрос задаёт $\ell=f(x_0)-G(x_0)$.
Если $f(x)-\ell\le\varepsilon$, останавливаемся.
Иначе берём $d=s-x$, начинаем с половины ранее принятой кривизны
(начальное сохранённое значение равно $L_{\rm init}$) и удваиваем
$L$, пока не выполнится
\begin{equation}\label{eq:stopping_bt}
 \alpha=\min\left\{1,\frac{G(x)}{L\|d\|^2}\right\},\qquad
 f(x+\alpha d)\le f(x)-\alpha G(x)+\frac{L\alpha^2\|d\|^2}{2}.
\end{equation}
Принимаем $x+\alpha d$, сохраняем значение функции и принятое $L$.
Это обычный подбор шага метода Франка-Вульфа \cite{pedregosa}.
Значения функции, градиенты и линейные минимумы на всём $Q$ точны.

\begin{theorem}[Конечная остановка и число всех запросов]
\label{thm:finite_stopping}
Каждый подбор шага завершается. Алгоритм останавливается с
$f(x_N)-f(x^*)\le f(x_N)-\ell_N\le\varepsilon$.
Если $G(x_0)>\varepsilon$, положим
\[
 J=\left\lceil\log_2\frac{G(x_0)}{\varepsilon}\right\rceil,
 \quad b_j=G(x_0)/2^j\ (0\le j<J),
\]
\[
 H_{\rm init}=\left\lceil
 \log_2\max\left\{1,\frac{L_{\rm init}}{2K_Q}\right\}\right\rceil.
\]
Число $N$ принятых шагов удовлетворяет
\begin{align}
 N&\le H_{\rm init}+\left\lceil\sum_{j=0}^{J-1}
       \max\left\{4,\frac{16K_QD^2}{b_j}\right\}\right\rceil
       \label{eq:stopping_steps}\\
  &\le H_{\rm init}+\left\lceil4J+
       \frac{32K_QD^2}{\varepsilon}\right\rceil.
       \label{eq:stopping_simple}
\end{align}
Вычислений градиента и запросов линейной минимизации ровно по $N+1$,
включая последний запрос. Начальное значение функции и все принятые
и отклонённые пробы при подборе шага дают не более
\begin{equation}\label{eq:stopping_values}
 N_f\le 1+2N+
 \left\lceil\log_2\max\left\{1,\frac{2K_Q}{L_{\rm init}}\right\}\right\rceil
\end{equation}
вычислений функции.
Если $G(x_0)\le\varepsilon$, шагов нет, а каждый из трёх оракулов
вызывается ровно один раз; величины для фаз не нужны.
\end{theorem}

\begin{proof}
Каждая сохранённая $\ell$ ограничивает $f(x^*)$ снизу, поэтому
\[
 0\le f(x)-f(x^*)\le f(x)-\ell\le G(x).
\]
В нетерминальной точке $G(x)>0$ и $d\ne0$.
Любое $L\ge K_Q$ проходит проверку \eqref{eq:stopping_bt}.
Для принятых значений выполнено
\begin{equation}\label{eq:stopping_progress}
 \begin{split}
 L_k&\le\max\{L_{k-1}/2,2K_Q\},\qquad L_{-1}=L_{\rm init},\\
 f(x_k)-f(x_{k+1})&\ge
 \min\left\{\frac{G(x_k)}2,\frac{G(x_k)^2}{2L_kD^2}\right\}.
 \end{split}
\end{equation}
Значение функции убывает, $\ell$ возрастает, а $f(x_k)-\ell_k$
убывает. Если $L_k>2K_Q$, была принята первая проба: иначе
предыдущая отклонённая проба $L_k/2>K_Q$ тоже прошла бы проверку.
Такой шаг делит сохранённую кривизну пополам. После достижения
$L_k\le2K_Q$ эта граница сохраняется. Шагов выше этого порога
не больше $H_{\rm init}$.

Каждый оставшийся нетерминальный шаг отнесём к фазе
$b_j/2<f(x_k)-\ell_k\le b_j$. Фазы посещаются по порядку.
В фазе $j$ зазор больше $b_j/2$, и \eqref{eq:stopping_progress}
даёт уменьшение функции не меньше
\[
 \min\left\{b_j/4,\ b_j^2/(16K_QD^2)\right\}.
\]
Суммарное уменьшение на всех шагах, начинающихся в этой фазе,
не больше истинной ошибки при входе в неё, то есть не больше $b_j$.
Сюда входит и последний шаг, выводящий из фазы: в его конечной
точке функция не меньше $f(x^*)$. Деление и суммирование доказывают
\eqref{eq:stopping_steps}. Из $\sum_{j<J}b_j^{-1}<2/\varepsilon$
следует \eqref{eq:stopping_simple}. Та же конечная граница верна
для каждого незавершённого начала процесса, поэтому бесконечный
процесс без остановки невозможен.

Если на шаге $k$ отклонены $q_k$ проб, то
$L_k=2^{q_k-1}L_{k-1}$. При $N\ge1$ получаем
\[
 \sum_{k<N}(q_k+1)=2N+\log_2\frac{L_{N-1}}{L_{\rm init}},
 \quad L_{N-1}\le\max\{L_{\rm init}/2^N,2K_Q\}.
\]
Добавление начального значения функции даёт
\eqref{eq:stopping_values}. В каждой посещённой точке нужен один
градиент и одно решение задачи линейной минимизации.
\end{proof}

\subsection{Перенос направленной точности между итерациями}
Направленный запрос с требуемой точностью $\eta>0$ возвращает
$h\ge0$ и $d=s-x$, $s\in Q$, с гарантированными границами
\begin{equation}\label{eq:directional_query}
 G(x)\le h+\eta,\qquad
 -\langle\nabla f(x),d\rangle\ge h-\eta.
\end{equation}
Это явное требование к оракулу. Проверка при подборе шага сама
по себе его не доказывает. Значения функции остаются точными.
Основной алгоритм с точными оракулами этого дополнения не требует.

Начнём с $\eta_0>0$ и зададим $0<\theta<1$.
После каждого запроса обновляем
$\ell\leftarrow\max\{\ell,f(x)-(h+\eta)\}$ и сначала проверяем
$f(x)-\ell\le\varepsilon$. Если остановки нет и $\eta>\theta h$,
делим $\eta$ пополам и повторяем запрос в той же точке, сохраняя
её значение функции. Когда $\eta\le\theta h$, используем
\eqref{eq:stopping_bt}, заменив $G(x)$ на $h-\eta$.
В следующую точку переносим достигнутую $\eta$.

\begin{corollary}[Учёт всех уточнений точности]
\label{cor:directional_stopping}
Правило останавливается с $f(x_N)-f(x^*)\le\varepsilon$.
Положим $\kappa=(1-\theta)/(1+\theta)$. Общее число $R$ делений
точности пополам удовлетворяет
\begin{equation}\label{eq:accuracy_refinements}
 R\le\left\lceil\log_2\max\left\{1,
 \frac{(1+\theta)\eta_0}{\theta\varepsilon}\right\}\right\rceil.
\end{equation}
Если первая верхняя граница зазора $h_0+\eta_0>\varepsilon$,
определим $J$ и $b_j$ как выше, заменив $G(x_0)$ на $h_0+\eta_0$.
Тогда
\begin{equation}\label{eq:directional_steps}
 N\le H_{\rm init}+\left\lceil
 \frac{4J}{\kappa}+\frac{32K_QD^2}{\kappa^2\varepsilon}
 \right\rceil.
\end{equation}
Направленных запросов ровно $N+1+R$. Граница
\eqref{eq:stopping_values} для вычислений функции сохраняется.
Если каждый направленный запрос использует один линейный минимум
на всём $Q$, таких вызовов тоже $N+1+R$.
Если первая граница зазора не больше $\varepsilon$, правило
останавливается сразу.
\end{corollary}
\begin{proof}
При неудачной проверке точности без остановки
\[
 \varepsilon<f(x)-\ell\le h+\eta
 <(1+\theta)\eta/\theta.
\]
Так как $\eta$ делится пополам и не увеличивается, получаем
\eqref{eq:accuracy_refinements}. При успешной проверке без остановки
\[
 h-\eta\ge\kappa(h+\eta)\ge\kappa(f(x)-\ell)>0.
\]
Предыдущее доказательство уменьшения и подсчёта шагов применимо
с порогом направленного спуска $\kappa b_j/2$, что даёт
\eqref{eq:directional_steps}. В каждой посещённой точке есть один
завершающий запрос, включая последнюю точку, а каждое деление
точности добавляет один неудачный запрос. При уточнении в одной
точке значение функции повторно не вычисляется.
\end{proof}

Доказательство остановки использует полный зазор и допускает
ненулевой предельный градиент. В разделе~\ref{sec:normal} и точная
граница кривизны, и всё выполнение метода не зависят от сколь угодно
большой нормальной компоненты градиента. Число направленных запросов
указывает, как часто нужна заданная точность; стоимость получения
каждого ответа определяется реализацией оракула.
